\pdfoutput=1
\documentclass[11pt]{article}
\usepackage[a4paper,margin=25mm]{geometry}
\usepackage{amsmath,amssymb,amsthm}
\usepackage[T1]{fontenc}
\usepackage{lmodern}
\usepackage[hidelinks]{hyperref}
\newtheorem{proposition}{Proposition}
\title{A counterexample to the minimum order clause\protect\\in TLMC 00000008420}
\author{Galaxy-0}
\date{5 October 2026}
\begin{document}
\maketitle
\begin{abstract}
The conjecture includes the assertion that the smallest order of a Kirkman
triple system admitting a transitive automorphism is $15$. Under its stated
definition, the one-block system of order $3$ is a counterexample. Its
three-cycle is a single transitive automorphism, including preservation of
the resolution. Thus the conjecture, as a conjunction, is false.
\end{abstract}

\section{The clause and its definitions}
The English statement defines a Kirkman triple system as a resolvable
Steiner triple system and claims that ``the smallest order of a KTS with a
transitive automorphism is the point-transitive type of order 15.''
We use the stronger literal interpretation: a single automorphism must
generate a point-transitive cyclic action.

A Steiner triple system on a finite set $V$ has a collection $\mathcal B$
of three-element subsets such that each pair of distinct points belongs
to exactly one block. A resolution is a partition $\mathcal R$ of
$\mathcal B$ into parallel classes, each of which partitions $V$.
An automorphism of the resolved system is a permutation of $V$ that
preserves the blocks and the resolution. The statement does not impose
$n>3$ or require a nontrivial system.

\section{The counterexample}
\begin{proposition}
There exists a Kirkman triple system of order $3$ with a single transitive
automorphism. Consequently the asserted minimum order $15$ is false.
\end{proposition}
\begin{proof}
Set
\[
 V=\{0,1,2\},\qquad B=\{0,1,2\},\qquad
 \mathcal B=\{B\},\qquad \mathcal R=\{\{B\}\}.
\]
The only block has three distinct points. Each of the three unordered
pairs of distinct points lies in $B$, and there is no second block.
Hence $(V,\mathcal B)$ is a Steiner triple system. The sole parallel
class $\{B\}$ partitions $V$, and the collection containing that class
partitions $\mathcal B$. This is a Kirkman triple system.

Let $\sigma=(0\ 1\ 2)$. It is a bijection, fixes $B$ setwise, and fixes
the sole parallel class. It is therefore an automorphism of the resolved
system. Its iterates take any given point to all three points, so the
cyclic group generated by this single automorphism acts transitively.
In particular, $\sigma$ has order exactly $3$.

The claimed minimum would imply that every such system has order at
least $15$. The constructed system has order $3<15$, a contradiction.
\end{proof}

\section{Scope and the earlier rejected argument}
This argument refutes the minimum-order conjunct. It makes no claim
against the existence criterion or the formula for the number of
parallel classes. Indeed, the example has $(3-1)/2=1$ parallel class.
It also does not refute the separate assertion about a full
automorphism group of prime cyclic order: the full automorphism group
of this example is $S_3$, not $C_3$. The relevant fact is that it
\emph{contains} the transitive three-cycle.

The earlier submission in pull request 21 used the order-$9$ affine
plane and a transitive translation group. The subsequent audit removed
that submission because a transitive group does not establish the
existence of a single transitive automorphism. A nonzero translation of
$\mathbb F_3^2$ has order $3$ and three orbits of size $3$, so it cannot
be transitive on nine points. The audit additionally observes that
$\operatorname{AGL}(2,3)$ has no element of order $9$. The present
counterexample avoids that substitution by displaying the required
single permutation explicitly. It addresses the statement as written;
an amended statement excluding order $3$ would require a different
analysis.

The order-$3$ convention is standard, not an invented relaxation.
Brown and Mellinger's discussion explicitly includes a KTS with three
points, one block, and one parallel class, while distinguishing the
smallest nontrivial case of order $9$ [2, PDF page 3].

\section{Formalization}
The accompanying \texttt{Counterexample.lean} imports only \texttt{Std}.
Its finite incidence model separately indexes points, blocks, and
resolution classes. It states and checks: exactly three distinct points
in each block; unique containment of each pair; unique containment of
each point within each resolution class; bijectivity on all three index
sets; preservation of incidence and resolution; and transitivity under
iterates of the same point permutation. It also checks that this
permutation has exact order $3$.

The theorem \texttt{minimum\_order\_clause\_false} negates a
minimum-order assertion with both attainment at $15$ and a lower-bound
condition. The final conjunction theorem records why falsifying this
conjunct suffices, without purporting to prove or disprove the other
conjuncts. There are no admitted proofs or added axioms in the source.

\begin{thebibliography}{9}
\bibitem{statement} The Last Math Competition,
\emph{Conjecture 00000008420}, repository statement.
\href{https://github.com/The-Last-Math-Competition/The-Last-Math-Competition/blob/d7d4bc9f6918401d3fb22c6e473df63eba5dbb78/conjectures/00000008420.md}{Repository statement}.
\bibitem{brown} E. Brown and K. E. Mellinger,
\emph{Kirkman's Schoolgirls Wearing Hats and Walking through Fields of Numbers},
Mathematics Magazine 82 (2009), no. 1.
\url{https://personal.math.vt.edu/brown/doc/kirkman.pdf}.
\bibitem{prior} The Last Math Competition, pull request 21 and subsequent audit
\href{https://github.com/The-Last-Math-Competition/The-Last-Math-Competition/commit/541cf4fbc7aef3e17085577f6403c11c09f31fe1}{audit commit 541cf4f}.
\href{https://github.com/The-Last-Math-Competition/The-Last-Math-Competition/pull/21}{Earlier pull request}.
\end{thebibliography}
\end{document}
